%MIT OpenCourseWare: https://ocw.mit.edu
%18.100A / 18.1001 Real Analysis, Fall 2020
%License: Creative Commons BY-NC-SA 
%For information about citing these materials or our Terms of Use, visit: https://ocw.mit.edu/terms.

%\subsection*{Lecture 4:}
%\textbf{\underline{The Characterization of the Real Numbers}}

\begin{question}
Last time we stated that $\QQ$ was an example of a field, but what \textit{is} a \textbf{field}?
\end{question}
\begin{definition}[Field]
A set $F$ is a \vocab{field} if it has two operations: addition $(+)$ and multiplication ($\cdot$) with the following properties.
\begin{enumerate}
    \item[A1)] If $x,y\in F$ then $x+y \in F$.
    \item[A2)] \textit{(Commutativity)} $\forall x,y\in F$, $x+y= y+x$.
    \item[A3)] \textit{(Associativity)} $\forall x,y,z\in F$, $(x+y)+z = x+(y+z)$.
    \item[A4)] $\exists$ an element $0\in F$ such that $0+x = x = x+0$.
    \item[A5)] $\forall x\in F$, $\exists y\in F$ such that $x+y=0$. We denote $y=-x$.
    \item[M1)] If $x,y\in F$, then $x\cdot y\in F$.
    \item[M2)] \textit{(Commutativity)} $\forall x,y\in F$, $x\cdot y= y\cdot x$.
    \item[M3)] \textit{(Associativity)} $\forall x,y,z\in F$, $(x\cdot y)\cdot z= x\cdot (y\cdot z)$.
    \item[M4)] $\exists 1\in F$ such that $1\cdot x = x = x\cdot 1$ for all $x\in F$.
    \item[M5)] $\forall x\in F\setminus\{0\}$, $\exists x^{-1}$ such that $x\cdot x^{-1} = 1.$
    \item[D)] \textit{(Distributativity)} $\forall x,y,z\in F$, $(x+y)z = xz+yz.$
\end{enumerate}
\end{definition}
These may seem like trivial properties, but consider the following non-example: $\ZZ$. $\ZZ$ fails M5)-- multiplicative inverses do not exist in the integers.

\begin{example}
Here are two examples of fields:
\begin{enumerate}
    \item $\ZZ_2 = \{0,1\}$ where $1+1 = 0$.
    \item $\ZZ_3 = \{0,1,2\}$ with $c:= a+b \pmod{3}$. In other words, \[2+1=3=0\pmod 3 \hspace{.25cm} \text{and} \hspace{.25cm} 2\cdot2 = 4 = 3+1 = 1\pmod 3.\]
\end{enumerate}
\end{example}

Simple properties follow from the properties of a field!
\begin{theorem}
If $x\in F$ where $F$ is a field, $0x=0.$
\end{theorem}
\textbf{Proof}: If $x\in F$, then
\[
0 = 0\cdot x -0\cdot x = (0+0)\cdot x -0\cdot x = 0\cdot x+0\cdot x -0\cdot x = 0\cdot x.
\]
\qed 

\begin{definition}[Ordered field]
A field $F$ is an \vocab{ordered field} if $F$ is also an ordered set with ordering $<$ and 
\begin{enumerate}[label = \roman*)]
    \item $\forall x,y,z\in F$, $x<y\implies x+z<y+z$.
    \item If $x>0$ and $y>0$ then $xy>0$.
\end{enumerate}
If $x>0$ we say $x$ is \textbf{positive}, and if $x\geq 0$ we say $x$ is \textbf{non-negative}.
\end{definition}

\begin{example}
$\QQ$ is an ordered field.
\end{example}

A non-example would be $\ZZ_2$. For instance, consider $0<1$. If $0<1$, then $1+0<1+1 = 0\implies 1<0$ which is a contradiction. If $1<0$, then $0=1+1<0+1 \implies 0<1$ which is a contradiction. Hence, $\ZZ_2$ is not an ordered field.

Using the definition of an ordered field, one can easily prove all of the usual facts about inequalities.
\begin{theorem}
If $x>0$, then $-x<0$ (and vice versa).
\end{theorem}
\textbf{Proof}: If $x\in F$ and $x>0$, then by i),
\[
-x+x>-x \implies 0>-x.
\]\qed 

One can see Proposition 1.1.8 \textbf{[L]} for a list of other simple inequality facts.

\begin{theorem}
Let $x,y\in F$ where $F$ is an ordered field. If $x>0$ and $y<0$ or $x<0$ and $y>0$, then $xy<0$.
\end{theorem}
\textbf{Proof}: Suppose $x>0$ and $y<0$. Then, $x>0$ and $-y>0$. Hence, $-xy = x(-y)>0$. Thus, $xy<0$. If $x<0$ and $y>0$, then $-x>0$ and $y>0\implies -xy = (-x)y>0 \implies xy<0.$ \qed

\begin{question}
Is there a greatest lower bound property? %how to word?
\end{question}
For an ordered field $F$, if $F$ has the least upper bound  property then $F$ has a greatest lower bound property.

\begin{theorem}
Let $F$ be an ordered field with the least upper bound property. If $A\subset F$ is nonempty and bounded below, then $\inf A$ exists in $F$.
\end{theorem}
\textbf{Proof}: Suppose $A\subset F$ is nonempty and bounded below, i.e. $\exists a\in F$ such that $\forall x\in A$, $a\leq x$. Let\\ $B = \{-x\mid x\in A\}$. Then, $\forall x\in A$, $-x\leq -a\implies -a$ is an upper bound for $B$. Since $F$ has the least upper bound property, $\exists c\in F$ such that $c = \sup B$. Then, $\forall x\in A$, $-x\leq c\implies \forall x\in A,$ $-c\leq x$. Hence, $-c$ is a lower bound for $A$. We have also shown that if $a$ is a lower bound for $A$, then $-a$ is an upper bound for $B$. Therefore, $c\leq -a$ since $c=\sup B\implies a\leq -c$. Hence, $-c$ is the greatest lower bound for $A$.
\qed 

\noindent \underline{\textbf{Real Numbers}}

\begin{theorem}
There exists a "unique" ordered field, labeled $\RR$, such that $\QQ\subset \RR$ and $\RR$ has the least upper bound property.
\end{theorem}
One can construct $\RR$ using Dedekind cuts or as equivalence classes of Cauchy sequences. (We will define Cauchy sequences later in the course.)

\begin{theorem}
$\exists! r\in \RR$ such that $r>0$ and $r^2 =2$. In other words, $\sqrt 2 \in \RR$ but $\sqrt 2 \notin \QQ$.
\end{theorem}
\textbf{Proof}: Let $\Tilde{E} = \{x\in \RR \mid x>0 \text{~~and~~} x^2<2\}$. Then, since $\Tilde{E}$ is bounded above (by 2 for instance), we have that $r:= \sup\Tilde{E}$ exists in $\RR$. Then, one can show that $r>0$ and $r^2 = 2$. This is left as an exercise.

We now prove uniqueness. Suppose that there is a $\Tilde{r}>0$ with $\Tilde{r}^2 = 2$. Then, since $(r+\Tilde{r})>0$,
\[
0 = r^2 -\Tilde{r}^2 = (r+\Tilde{r})(r-\Tilde{r}) \implies r-\tilde r = 0 \implies r = \Tilde{r}.
\] 
\qed 

\begin{remark}
In Assignment 2 Exercise 7, you will show that $\sqrt[3]2\in \RR$.
\end{remark}